\SourceLeaf{358}%
gesamten Erdkörpers immer nur kleine Größen sein. Das erlaubt Umformungen, durch die die Formeln erheblich flexibler und bequemer werden. Obwohl die umgeformten Ausdrücke eigentlich nur Näherungsformeln sind, können sie sogar dieselbe oder eine noch höhere Genauigkeit erreichen als die absolut strengen Formeln. Das wird man. nicht paradox finden, wenn man bedenkt, dass diese doch immer mithilfe trigonometrischer Tafeln angewandt werden müssen, deren Genauigkeit endlich und durch die Zahl der Dezimalstellen begrenzt ist. Unter den verschiedenen Formeln, die in der ersten Abhandlung für den genannten Zweck angegeben wurden, zeichnet sich besonders die Kombination am Schluss aus. Sie stellt den Zusammenhang zwischen jenen sechs Größen in der für die Rechnung bequemsten Form dar und bietet selbst bei den größten tatsächlich messbaren Dreiecken mehr als genug Genauigkeit. Dadurch musste der Wunsch nach entsprechenden Formeln entstehen, die unmittelbar für die Ellipsoidoberfläche gelten. Ihre Entwicklung bildet den Hauptinhalt der vorliegenden \textit{zweiten} Abhandlung.

Die genannten Formeln für die Kugeloberfläche ließen sich aus ganz elementaren Sätzen gewinnen. Ihre Gegenstücke für die Ellipsoidoberfläche zu ermitteln, erfordert dagegen eine Reihe ziemlich verwickelter Rechenschritte. Daher ist es zweifellos willkommen, wenn mehr als ein Weg zum selben Ziel aufgezeigt wird. Der Verfasser hatte alle diese Untersuchungen bereits vor mehr als dreißig Jahren für seinen eigenen Gebrauch durchgeführt, aber bisher keinen besonderen Anlass zur Veröffentlichung gefunden. In der vorliegenden Abhandlung teilt er nun \textit{zwei} völlig verschiedene Herleitungen mit, die am Ende zu genau denselben Ergebnissen führen. Eine davon beruht auf der Theorie der konformen Abbildung der Ellipsoidoberfläche auf die Kugeloberfläche. In dieser Hinsicht schließt die zweite Abhandlung an die erste an. Im Übrigen sind beide insofern als ganz unabhängig voneinander anzusehen, als es freisteht, die geodätischen Rechnungen ausschließlich nach der Methode der ersten oder ausschließlich nach der Methode der zweiten Abhandlung zu führen.

Die Aufgabe, von der geografischen Lage eines Punkts auf der Sphäroidoberfläche zu der eines anderen Punkts überzugehen, der mit ihm durch eine %
